\section{A bounded-metric criterion for klt pairs}\label{sec:criterion}

Throughout this section $(W,D)$ satisfies the hypotheses of Theorem~\ref{thm:klt}, and $P=-(K_W+D)$. Divisor equalities involving line bundles mean linear equivalence; the corresponding equalities for rational divisors mean $\QQ$-linear equivalence. An assertion about an effective divisor, its support, or its coefficient is an assertion about the actual divisor. We will explicitly distinguish these uses.

\subsection{Analytic inputs and numerical comparisons}\label{subsec:inputs}

We record the analytic results in the forms used below. If $B$ is a line bundle on a compact K\"ahler $n$-fold $S$ with a semipositive singular metric $h$, the hard Lefschetz theorem \cite[Theorem~2.1.1]{DPS2001} gives a surjection
\begin{equation}\label{eq:HL}
H^0(S,\Omega_S^{n-q}\otimes B\otimes\II(h))
\longrightarrow H^q(S,K_S\otimes B\otimes\II(h)).
\end{equation}
Here $\II(h)$ tests local integrability of $|g|^2e^{-\varphi}$ in a frame of weight $\varphi$. In particular all positive powers of a bounded metric have trivial multiplier ideal.

For a projective surjective morphism $f:S\to Y$ of smooth varieties with connected fibers, the singular adjoint direct-image theorem \cite[\S3.2.1 and Theorems~3.3.4--3.3.5]{PaunTakayama2018} applies to a semipositive singular metric on $B$. When its multiplier ideal is trivial on the total space and on general fibers, its fiberwise $L^2$ metric on $f_*(K_{S/Y}+B)$ has semipositive curvature in the singular sense. On a base-change open where the direct image has rank one and its generator has finite integral, the negative logarithm of that integral is a psh local weight. The integral formula and the metric can be compared almost everywhere; we do not need continuity or base change at every point. We apply this statement separately for every integer twist, so the base-change open may depend on that integer.

Finally, the theorem of Fujino \cite[Theorem~1.3]{Fujino2018} states that if $f:S\to Y$ is a surjective morphism from a compact K\"ahler manifold to a projective variety, and a singular metric on a line bundle $B$ has curvature at least $\epsilon f^*\omega_A$ for $\epsilon>0$ and a positive curvature form of an ample line bundle $A$ on $Y$, then
\begin{equation}\label{eq:KN}
H^i\bigl(Y,R^j f_*(K_S\otimes B\otimes\II(h))\bigr)=0
\quad (i>0,\ j\ge0).
\end{equation}
Only $j=0$ will be needed.

We also need a numerical fact that explains the role of $\QQ$-factoriality. For a birational morphism $\pi:S\to W$ with $S$ smooth and $W$ normal and $\QQ$-factorial, pushforward of divisors induces a linear map on numerical divisor classes and preserves pseudoeffectivity; compare \cite[Lemma~3.4]{LMPTX2023}. Indeed, if a divisor $F$ on $S$ is numerically trivial, then $\pi^*\pi_*F-F$ is exceptional and is numerically trivial over $W$. Applying the negativity lemma to this divisor and its negative makes it zero. Intersecting $\pi^*\pi_*F=F$ with curves dominating curves of $W$ proves that $\pi_*F$ is numerically trivial. Pushforward is consequently a well-defined continuous linear map on the finite-dimensional spaces of numerical divisor classes. It sends effective divisors to effective divisors, hence preserves their closed cone. Pullback of a pseudoeffective $\QQ$-Cartier class by the birational morphisms used here also preserves pseudoeffectivity.

\subsection{Twisted differential forms and determinants}\label{subsec:determinant}

The first task is to produce an unbounded sequence of effective divisors differing by multiples of $P$ from one fixed class. We adapt the determinant argument of Lazi\'c--Peternell \cite[Lemma~4.1]{LP2018}, as used in \cite[Lemma~5.1]{LMPTX2023}. Here the original twisting exponents are positive and unbounded, so a basis containing a section of high exponent has an unbounded sum of exponents. This direct observation is sufficient in our setting; we do not import the numerical-nontriviality hypothesis of the cited lemma.

\begin{lemma}\label{lem:determinant}
Either $H^0(W,mP)\ne0$ for some positive Cartier multiple, or there are a Weil divisor $M$, strictly increasing positive integers $m_i$, and effective Weil divisors $N_i$ such that
\begin{equation}\label{eq:Ni}
N_i\sim_{\QQ} M+m_iP,\qquad -M\in\PE(W).
\end{equation}
All the $m_iP$ may be required to be Cartier.
\end{lemma}

\begin{proof}
The assertion is immediate if $W$ is a point. Suppose $\dim W=n>0$ and choose a smooth projective resolution $\rho:\widetilde W\to W$ on which the metric is bounded. Let $r_0>0$ make $r_0P$ Cartier. Riemann--Roch makes
\[
t\longmapsto\chi(\widetilde W,K_{\widetilde W}+t\rho^*(r_0P))
\]
a polynomial. Its value at $t=0$ is $(-1)^n\chi(\widetilde W,\OO_{\widetilde W})$, which is nonzero by hypothesis. It is therefore nonzero for all but finitely many positive integers. For each such integer at least one cohomology group is nonzero. Passing to an infinite subsequence fixes a cohomological degree $q$. Apply \eqref{eq:HL} with $B=t\rho^*(r_0P)$. Its multiplier ideal is trivial. We obtain a fixed $p=n-q$ and nonzero sections
\[
\sigma_i\in H^0(\widetilde W,\Omega^p_{\widetilde W}\otimes\OO_{\widetilde W}(n_i\rho^*P)),
\qquad n_i\longrightarrow\infty.
\]
If $p=0$, projection to $W$ gives the first alternative because $\rho_*\OO_{\widetilde W}=\OO_W$.

Assume $p>0$. Work over the function field $K=\CC(W)$ and choose a rational trivialization for $r_0P$. Untwisting the sections defines vectors in $\bigwedge^p\Omega^1_{K/\CC}$. Let $F_K$ be their span and $b=\dim_KF_K$. Its determinant is a one-dimensional subspace of
\[
\bigwedge^b\bigl(\bigwedge^p\Omega^1_{K/\CC}\bigr)
\subset (\Omega^1_{K/\CC})^{\otimes pb}.
\]
The indicated inclusion uses alternating tensors in characteristic zero. Intersect this line with $(\Omega_W^1)^{[\otimes pb]}$, the reflexive tensor power, and take its saturated reflexive rank-one extension. This is a sheaf $\OO_W(M)$ for a Weil divisor $M$. The determinant of any basis selected from the $\sigma_i$ gives a nonzero section of $\OO_W(M+\sum n_iP)$ on the smooth locus of $W$. It is regular at every codimension-one point: the resolution is an isomorphism there, and the original sections are holomorphic. Reflexivity extends it across codimension two. Its zero divisor is an effective Weil divisor $N$.

The sums of exponents of such bases are unbounded. Given any one of the nonzero vectors at arbitrarily high exponent, extend it to a basis using vectors from the spanning family. All exponents are positive, so the sum is at least the chosen exponent. Selecting a strictly increasing subsequence of these sums gives $m_i$ and $N_i$ in \eqref{eq:Ni}.

The pair $(W,D)$ is projective and klt, the boundary is effective, and $-(K_W+D)$ is nef. Thus \cite[Theorem~4.1]{LMPTX2023} applies to the saturated rank-one subsheaf $\OO_W(M)$ of the reflexive cotangent tensor power. It asserts that its negative first Chern class is pseudoeffective. Since $W$ is $\QQ$-factorial, this is the claimed pseudoeffectivity of the $\QQ$-Cartier divisor $-M$.
\end{proof}

\subsection{A maximal dominated fibration}\label{subsec:maximal}

Suppose we are in the second case of Lemma~\ref{lem:determinant}. Our next task is to express a multiple of a birational pullback of $P$ as the pullback of a line bundle from a suitable base plus an effective divisor. We need a further property: above each prime divisor of the base, some component must be absent from that effective divisor. This property will prevent the induced base metric from acquiring a divisorial singularity. To obtain the decomposition, we choose a base of maximal dimension whose polarization is dominated by $P$ in pseudoeffective order. Maximality will exclude pencils satisfying the corresponding domination condition that are nonconstant on a general fiber.

Consider diagrams
\begin{equation}\label{eq:diagram}
\begin{tikzcd}[column sep=large]
S \arrow[r,"f"] \arrow[d,"\pi"'] & Y\\
W &
\end{tikzcd}
\end{equation}
where $S,Y$ are smooth projective varieties, $\pi$ is birational, $f$ is surjective, and
\begin{equation}\label{eq:domination}
d\pi^*P-f^*A\in\PE(S)
\end{equation}
for an integer $d>0$ and an ample line bundle $A$ on $Y$. The map to a point is included, with the polarization omitted. Choose such a diagram with $\dim Y$ maximal.

We can require connected general fibers. Passing to the relative algebraic closure of $\CC(Y)$ in $\CC(W)$ gives a generically finite replacement of the base. The pullback of $A$ there is big. On a smooth projective model it dominates a small ample class in pseudoeffective order, so after multiplying \eqref{eq:domination} its required form persists. The same reasoning shows that smooth birational modifications of the base are harmless. In particular $\CC(Y)$ is algebraically closed in $\CC(W)$ for the resulting fibration. The finite map in the Stein factorization of the proper morphism $f$ is then birational. Since $Y$ is normal, that map is an isomorphism. Thus every fiber of $f$ is connected and $f_*\OO_S=\OO_Y$, as required by the direct-image theorem.

There is also a normal equidimensional intermediate model
\begin{equation}\label{eq:intermediate}
\begin{tikzcd}[row sep=large,column sep=large]
& S \arrow[d] \arrow[ddl,bend right=25,"f"'] \arrow[ddr,bend left=25,"\pi"] &\\
& S_0 \arrow[dl,"f_0"'] \arrow[dr,"\pi_0"] &\\
Y && W .
\end{tikzcd}
\end{equation}
Here $\pi_0$ is a birational morphism and $S$ resolves $S_0$ and the pair on $W$. To construct it, start with a projective graph of the rational map, flatten its main component by a projective modification of the base \cite{RaynaudGruson1971}, replace the base by a smooth projective modification, and normalize the main flat transform. Finite normalization preserves equidimensionality over the smooth base. Finally resolve above this model, also dominating the resolution on which the given metric is bounded. The smooth space $S$ need not be equidimensional over $Y$. We use equidimensionality only on $S_0$, where it implies that every vertical prime divisor dominates a prime divisor of $Y$.

\begin{lemma}[Excluding a new pencil]\label{lem:pencil}
In a maximal diagram \eqref{eq:diagram}, let $B$ be a line bundle on a birational model over $S$. Suppose that its class is pseudoeffectively dominated by $c\pi^*P+E'$ for some $c>0$ and a $\pi$-exceptional divisor $E'$. Then no pencil in $|B|$ is nonconstant on the general fiber of $f$. The conclusion still applies after adding to $B$ the pullback of any fixed divisor on $Y$.
\end{lemma}

\begin{proof}
Resolve the pencil, remove its fixed part, and combine its morphism to $\mathbb P^1$ with $f$. If the pencil is nonconstant on the general fiber, the image of the combined map has dimension greater than $\dim Y$. A product polarization gives a basepoint-free line bundle $F$ on the resolved source, dominated by a multiple of $\pi^*P$ plus an exceptional divisor, using \eqref{eq:domination}.

Choose a general member $F_0$ of that system. It has no $\pi$-exceptional component. Push the class comparison down to $W$, using the numerical fact in Section~\ref{subsec:inputs}; the exceptional term vanishes. The result says that $c'P-\pi_*F_0$ is pseudoeffective. Pulling back gives $c'\pi^*P-\pi^*\pi_*F_0$ pseudoeffective. Since the effective divisor $\pi_*F_0$ is $\QQ$-Cartier, its pullback contains its strict transform $F_0$ and has an effective exceptional remainder. Hence $c'\pi^*P-F_0$ is pseudoeffective. This dominates a polarization for the larger projective image. On a smooth model of that image the pulled-back polarization is big, so it dominates a small ample class; scaling and pulling back gives a diagram of the allowed form. This contradicts maximality.

For the last assertion, choose an integer $b$ such that $bA$ minus the fixed base divisor is ample. The extra pullback is then dominated by $bf^*A$, hence by $bd\pi^*P$ using \eqref{eq:domination}. Increase $c$ accordingly.
\end{proof}

\subsection{Horizontal affineness and vertical minima}\label{subsec:minima}

Pull the effective $\QQ$-Cartier divisors $N_i$ to $S_0$. We next produce the decomposition that will control a metric at the boundary of the base.

\begin{lemma}\label{lem:decomposition}
For some integer $r>0$ with $rP$ Cartier, there are a line bundle $V$ on $Y$ and an effective Cartier divisor $H$ on $S$ such that
\begin{equation}\label{eq:decomposition}
r\pi^*P=f^*V+H.
\end{equation}
For each prime divisor $G$ of $Y$, at least one component of $f^*G$ dominating $G$ has coefficient zero in $H$.
\end{lemma}

\begin{proof}
For $m_i<m_j<m_k$, set
\[
\alpha=\frac{m_k-m_j}{m_k-m_i},\qquad
\beta=\frac{m_j-m_i}{m_k-m_i}.
\]
Then $\alpha,\beta>0$, $\alpha+\beta=1$, and
\[
N_j\sim_{\QQ}\alpha N_i+\beta N_k.
\]
After clearing denominators, the two effective divisors give sections of the same line bundle. That bundle is pseudoeffectively dominated by a multiple of $P$, because $-M$ is pseudoeffective in \eqref{eq:Ni}. Lemma~\ref{lem:pencil} says that their ratio is constant on the general fiber. The ratio belongs to $\CC(Y)$, using relative algebraic closure. Its principal divisor therefore has no horizontal component on $S_0$.

It follows that the horizontal parts of $\pi_0^*N_i$ are affine functions of $m_i$ as actual divisors. Fixing the first two indices expresses them as $A_0+m_iB_0$, where the supports of $A_0,B_0$ are contained in one finite union of primes. Effectivity for an unbounded sequence of $m_i$ gives $B_0\ge0$. Consequently the difference of two members, divided by the difference of their exponents, is $\QQ$-linearly equivalent to $\pi_0^*P$ and has effective horizontal part. Clearing denominators gives a rational section of a positive Cartier multiple of $\pi_0^*P$ with no horizontal poles. Write its divisor as $T_0$.

For every base prime $G$, let
\[
v_G=\min_{F\mapsto G}
\frac{\operatorname{coeff}_F(T_0)}{\operatorname{coeff}_F(f_0^*G)}.
\]
The denominator is positive, and the minimum is over the finitely many primes dominating $G$. Only finitely many $v_G$ are nonzero. Every vertical prime of $S_0$ occurs above a base prime, by equidimensionality. Thus
\[
T_0-f_0^*\Bigl(\sum_Gv_GG\Bigr)\ge0.
\]
The difference is $\QQ$-Cartier: $T_0$ is Cartier up to its chosen multiple, and $Y$ is smooth. Multiply again to make the base divisor integral and the remainder Cartier. This gives $V$ and an effective Cartier remainder on $S_0$. Pull them to $S$ to obtain \eqref{eq:decomposition}. The component attaining the minimum has zero remainder coefficient; its strict transform on $S$ still dominates $G$ and has coefficient zero in $H$.
\end{proof}

If $Y$ is a point, \eqref{eq:decomposition} immediately gives an effective divisor linearly equivalent to $r\pi^*P$, and hence a section downstairs. From now on suppose $\dim Y>0$. We will construct a semipositive metric on $V$ from fiber integrals and use the component absent from $H$ over each base prime to prove that this metric is bounded. The integrals must come from adjoint line bundles whose fiberwise section spaces have dimension one. Discrepancy rounding supplies these line bundles, and the same maximality argument supplies their rank-one property.

\subsection{Discrepancy rounding and rank one for every twist}\label{subsec:rankone}

Use compatible canonical divisors and put
\begin{equation}\label{eq:rounding}
R=K_S-\pi^*(K_W+D),\qquad E=\lceil R\rceil,
\qquad\Xi=\lceil R\rceil-R.
\end{equation}
Choose the resolution so that the support of $R$ is simple normal crossings. Since the pair is klt, all coefficients of $R$ are greater than $-1$. Along a nonexceptional prime the coefficient is minus its coefficient in $D$, which lies in $(-1,0]$; its ceiling is zero. Thus $E$ is integral, effective and $\pi$-exceptional. The divisor $\Xi$ is effective with simple normal crossings support and all coefficients in $[0,1)$. The identity we will use is
\begin{equation}\label{eq:adjointidentity}
K_S+\pi^*P+\Xi=E.
\end{equation}

\begin{lemma}\label{lem:rankone}
For every integer $a\ge0$,
\[
\rank f_*\OO_S(E+ar\pi^*P)=1.
\]
\end{lemma}
\begin{proof}
Over the generic fiber the section of $E+aH$ gives a nonzero section, by \eqref{eq:decomposition}. This includes $a=0$. If the dimension of the generic-fiber section space were at least two, twisting by the pullback of a sufficiently ample line bundle on $Y$ would extend two independent generic sections to global sections. Their ratio would be nonconstant on the general fiber. Their line bundle is dominated by a multiple of $\pi^*P$ plus the exceptional divisor $E$, including the base twist as in Lemma~\ref{lem:pencil}. This contradicts that lemma.
\end{proof}

For $a\ge0$ define the integral line bundle
\begin{equation}\label{eq:Ba}
B_a=(ar+1)\pi^*P+\Xi=E+ar\pi^*P-K_S.
\end{equation}
The second expression proves integrality even though the two summands in the first expression can be rational divisors. The role of $\Xi$ is now precise: its fractional poles make the adjoint line integral, while keeping the multiplier ideal trivial. Initially equip it with the metric whose weights are $(ar+1)\varphi_P+\varphi_\Xi$. Here $\varphi_P$ denotes the bounded weights on $\pi^*P$, and $\varphi_\Xi$ is the divisorial weight. Its curvature is semipositive. Its multiplier ideal is trivial because the weights of $P$ are bounded and the coefficients of the simple normal crossings divisor $\Xi$ are less than one. The same is true on general fibers, after shrinking to where the boundary is relatively normal crossings.

\subsection{A bounded metric from normalized fiber integrals}\label{subsec:metric}

Our next goal is to put a bounded semipositive metric on the line bundle $V$ in \eqref{eq:decomposition}. This is where the component absent from $H$ over each base prime is used.

\begin{proposition}\label{prop:boundedbase}
The line bundle $V$ admits a singular Hermitian metric whose local weights are psh and locally bounded on all of $Y$.
\end{proposition}

\begin{proof}
Fix a dense Zariski open $Y^\circ$ over which $f$ is smooth with connected fibers, $\Xi$ is relatively simple normal crossings, and no fiber is contained in $E\cup H$. In a coordinate neighborhood in $Y^\circ$, choose frames $v$ of $V$ and $\eta$ of $-K_Y$. The section
\[
s=f^*v\,s_H
\]
of $r\pi^*P$ has divisor $H$ over that neighborhood. The section $s_Es^af^*\eta$ belongs to $K_{S/Y}+B_a$. By Lemma~\ref{lem:rankone}, it generates the direct image at a general point. On its base-change open, this generator locally identifies the adjoint direct image with $aV-K_Y$. Because the direct image has rank one, its semipositive $L^2$ metric has a psh weight equal to the negative logarithm of the squared $L^2$ norm of this generator. Dividing by $a$ gives the functions below. This is a local identification where the generator is a frame, and asserts no global invertibility at omitted base points. We obtain psh functions
\begin{equation}\label{eq:ua}
u_a(y)=-\frac1a\log\int_{S_y}t(x)^a\,d\mu_y(x),
\qquad t=|s|_{h_P^r}^2,\qquad a\ge1.
\end{equation}
The measure $\mu_y$ is the squared adjoint norm density of $s_Ef^*\eta$, using the metric on $\pi^*P+\Xi$. Locally it has a bounded positive factor, the squared modulus of a nonzero holomorphic section, and the factors $|z_j|^{-2c_j}$ with $c_j<1$. It is therefore finite on a general compact fiber and has the same null sets as a smooth volume measure. Zeros of $s_E$ and the divisorial poles occupy sets of measure zero. This description uses boundedness, not continuity, of $\varphi_P$.

For each $y_0\in Y^\circ$, choose a point on $S_{y_0}$ off $E\cup H\cup\Supp\Xi$. On a small product neighborhood for the smooth map, the holomorphic factors are bounded away from zero and the metric weights are bounded. There are a neighborhood of $y_0$, constants $c,b>0$, and open pieces in its fibers on which $t\ge c$ and whose $\mu_y$-mass is at least $b$. Thus
\[
\int_{S_y}t^a\,d\mu_y\ge b c^a,
\qquad
u_a(y)\le-\log c-\frac{\log b}{a}.
\]
These estimates extend the psh representatives across each exponent-dependent base-change exceptional set and bound them locally uniformly above on $Y^\circ$.

Take the countable intersection $Y^*$ of the full-measure sets where these integral formulas agree with the psh representatives for every positive integer $a$. This is a common full-measure subset of $Y^\circ$, not a claimed common Zariski open. For $y\in Y^*$, the elementary $L^a$-norm limit on the finite measure space $(S_y,\mu_y)$ gives
\begin{equation}\label{eq:esssup}
\lim_{a\to\infty}u_a(y)
=-\log\esssup_{S_y}|s|_{h_P^r}^2=:u(y).
\end{equation}
For completeness, if $M_y=\esssup t$, the upper bound is
$(\int t^a\,d\mu_y)^{1/a}\le M_y\mu_y(S_y)^{1/a}$.
For any $c<M_y$, the set $\{t>c\}$ has positive measure, giving the reverse limiting bound $\liminf(\int t^a\,d\mu_y)^{1/a}\ge c$. Let $c\uparrow M_y$.

The psh functions $u_a$ have locally uniformly bounded upper bounds. They cannot converge to $-\infty$ locally: on each relatively compact base neighborhood, properness, holomorphicity of $s$, and boundedness of $\varphi_P$ bound $t$ above. The measures $\mu_y(S_y)$ are bounded above there: in relative simple normal crossings charts their densities are at most $C\prod_j|z_j|^{-2c_j}$ with $c_j<1$, and properness gives a finite cover above a compact base neighborhood. Hence \eqref{eq:ua} also has a uniform lower bound on compact subsets of $Y^\circ$. Local compactness for psh functions now gives subsequential $L^1_{\mathrm{loc}}$ convergence, and \eqref{eq:esssup} identifies every subsequential limit almost everywhere. Consequently $u$ has a psh representative.

Under a frame change $v'=gv$, the function $t$ is multiplied by $|g|^2$ and $u$ changes to $u-\log|g|^2$. Changing $\eta$ contributes a term divided by $a$, which disappears in the limit. Thus $u$ defines the weights of a metric on $V$ over $Y^\circ$.

The normalized integrals have now produced positivity on the smooth part of the base. It remains to prove boundedness at the boundary; positivity alone would not suffice for the later multiplier-ideal calculation. On the inverse image of a relatively compact base coordinate neighborhood, properness and bounded weights of $h_P$ give a uniform upper bound for $|s|^2$. Formula~\eqref{eq:esssup} therefore gives a lower bound for $u$ even near a missing base fiber. For an omitted prime $G\subset Y$, choose the component $F$ above $G$ on which $H$ has coefficient zero. At a general point of $F$, the section $s$ is nonzero, the map $F\to G$ is submersive, and local transverse coordinates have the form
\[
z_1\circ f=w_1^e\,a(w),\qquad a(w)\ne0.
\]
Together with coordinates along $G$, this shows that every nearby general fiber meets a neighborhood of that point in an open set. On a smaller neighborhood $|s|_{h_P^r}^2\ge c>0$ because the holomorphic section is nonzero and the metric is bounded. The neighborhoods intersecting such a general fiber have positive $\mu_y$-measure, since that measure has the same null sets as smooth volume. No uniform positive lower bound on their measures is needed for an essential supremum. Thus the component may lie in $E$ or $\Xi$ without affecting this argument: on nearby general fibers their divisorial supports have measure zero. The essential supremum is at least $c$. Thus $u\le-\log c$ near a dense open subset of $G$.

The psh removable-singularity theorem extends $u$ across these codimension-one sets, using the upper bound just proved. What remains is an analytic subset of codimension at least two in the smooth base. Psh Hartogs extension extends $u$ there as well. Its lower bound persists by the upper-semicontinuous extension, and its psh representative has a local upper bound. The local extensions are unique and respect the frame transformations, giving the asserted bounded metric on $V$ over all of $Y$. This conclusion is local boundedness. We have used no continuity assertion for the original metric or for its essential supremum.
\end{proof}

\subsection{One small mixture and Euler interpolation}\label{subsec:interpolation}

The bounded metric on $V$ lets us represent every nonnegative integral twist $ar\pi^*P=af^*V+aH$ by a semipositive metric whose only unbounded weights come from the integral divisor $aH$. These poles will give precisely the ideal $\OO_S(-aH)$, canceling $aH$ in the adjoint bundle. To apply vanishing, we still need strict positivity in base directions. We put it on the fixed summand $\pi^*P+\Xi$, so the same construction also works for the zeroth twist.

\begin{lemma}\label{lem:mixture}
There is a metric on $\pi^*P+\Xi$ with curvature at least $\epsilon f^*\omega_A$ for some $\epsilon>0$ and with trivial multiplier ideal.
\end{lemma}
\begin{proof}
By \eqref{eq:domination}, $d\pi^*P-f^*A$ has a positive-current metric. Tensoring with a positive metric on $f^*A$ and taking a $d$th root gives a metric on $\pi^*P$ with local weights $\psi$ and curvature at least $d^{-1}f^*\omega_A$. Choose $0<\tau<1$ and mix these weights with the bounded metric:
\[
\varphi_\tau=(1-\tau)\varphi_P+\tau\psi.
\]
Its curvature is at least $(\tau/d)f^*\omega_A$. We show that $\varphi_\tau+\varphi_\Xi$ has trivial multiplier ideal for one sufficiently small $\tau$.

Choose $p>1$ so close to one that $pc_j<1$ for every coefficient $c_j$ of $\Xi$, and let $q=p/(p-1)$. On a log-resolution chart,
\[
e^{-p\varphi_\Xi}\asymp\prod_j|z_j|^{-2pc_j}
\]
is integrable. Skoda's local small-exponent integrability theorem \cite{Skoda1972} implies that $e^{-q\tau\psi}$ is integrable for sufficiently small $\tau$. A finite cover of the compact $S$ permits one choice of $\tau$ for all charts. H\"older's inequality makes $e^{-(\varphi_\Xi+\tau\psi)}$ integrable; the remaining bounded factor $e^{-(1-\tau)\varphi_P}$ does not affect integrability. This is exactly the triviality of the multiplier ideal. No statement about analytic singularities of $\psi$ is needed.
\end{proof}

\begin{proof}[Proof of Theorem~\ref{thm:klt}]
The earlier arguments already dispose of $\dim W=0$, the $p=0$ alternative of Lemma~\ref{lem:determinant}, and a zero-dimensional maximal base. Otherwise take $V,H,E,\Xi$ as constructed above. Let $h_V$ be the bounded semipositive metric on $V$ from Proposition~\ref{prop:boundedbase}. For every integer $a\ge0$, equip $B_a$ in \eqref{eq:Ba} with the tensor product of the metric from Lemma~\ref{lem:mixture}, the pullback of $h_V^{\otimes a}$, and the divisor metric of $aH$. More explicitly, if $g_H$ is a local equation of $H$, this metric has local weight
\[
\varphi_{\mathrm{mix}}+a f^*\varphi_V+a\log|g_H|^2,
\]
where $\varphi_{\mathrm{mix}}$ is a weight of the fixed metric from Lemma~\ref{lem:mixture} and $\varphi_V$ is a weight of $h_V$. Curvature is at least the same positive multiple of $f^*\omega_A$. Since $\varphi_V$ is bounded and $aH$ is integral, the multiplier ideal is exactly $\OO_S(-aH)$. Indeed, at the generic point of a component of $H$, the remaining psh weight is bounded above, so integrability forces divisibility by $g_H^a$. The quotient extends holomorphically across codimension two by normality. Conversely, after this division, the test is the trivial-ideal test in Lemma~\ref{lem:mixture}. Bounded weights do not change that ideal.

The adjoint identity becomes
\[
K_S+B_a-aH=E+af^*V.
\]
Apply \eqref{eq:KN} and the projection formula. Writing $\mathcal G=f_*\OO_S(E)$, we obtain
\begin{equation}\label{eq:allvanishing}
H^i(Y,\mathcal G\otimes\OO_Y(aV))=0
\quad\text{for every }i>0\text{ and every }a\ge0.
\end{equation}
The sheaf $\mathcal G$ is coherent and has rank one. Riemann--Roch on the smooth projective base makes
\[
Q(a)=\chi(Y,\mathcal G\otimes\OO_Y(aV))
\]
a polynomial in $a$. The effective divisor $E$ supplies a nonzero section of $\mathcal G$. At $a=0$, \eqref{eq:allvanishing} therefore gives $Q(0)=h^0(Y,\mathcal G)>0$. The polynomial is nonzero, so there exists a positive integer $a$ with $Q(a)\ne0$. Again by \eqref{eq:allvanishing}, $Q(a)=h^0(Y,\mathcal G\otimes\OO_Y(aV))$, so this space of sections is nonzero. Thus
\[
H^0(S,E+af^*V)\ne0.
\]
Multiplication by $s_H^a$ gives a nonzero section of $E+ar\pi^*P$. Normality and the effective exceptional nature of $E$ imply $\pi_*\OO_S(E)=\OO_W$: a rational function with possible poles only over codimension at least two is regular on $W$. The projection formula now gives a nonzero section of $arP$ on $W$. This is a positive Cartier multiple, as required.
\end{proof}
